\chapter{Vision du Projet}

\section{Problématique}

Les fraudes bancaires sont rares par rapport au volume total des transactions, mais leur impact financier, opérationnel et réputationnel est élevé. Leur faible fréquence rend la détection difficile: un modèle naïf peut obtenir une accuracy très élevée tout en manquant la majorité des fraudes.

Le projet répond à cette problématique par un produit IA capable d'identifier automatiquement les transactions suspectes et de fournir un score de risque exploitable par des analystes fraude.

\section{Objectifs}

Les objectifs opérationnels sont les suivants:

\begin{itemize}
    \item construire une chaîne reproductible depuis les données Kaggle jusqu'à un modèle servi par API;
    \item traiter explicitement le déséquilibre de classes;
    \item éviter les fuites de données lors du preprocessing, du sampling et de l'évaluation;
    \item exposer le modèle via une API REST et une interface utilisateur;
    \item suivre les expériences et les artefacts avec MLflow;
    \item orchestrer les étapes DataOps/MLOps avec Dagster;
    \item rendre le projet testable et déployable avec Docker et CI/CD.
\end{itemize}

\section{Utilisateurs Cibles}

\begin{table}[htbp]
\centering
\caption{Utilisateurs cibles et besoins}
\begin{tabularx}{\textwidth}{p{0.28\textwidth}X}
\toprule
Utilisateur & Besoin principal\\
\midrule
Analyste fraude & Prioriser les transactions risquées et consulter le score de fraude.\\
Equipe risque bancaire & Suivre les indicateurs de fraude, de volumétrie et de performance.\\
Equipe Data/ML & Maintenir les pipelines, reproduire l'entraînement et contrôler le modèle.\\
Exploitant applicatif & Vérifier la disponibilité, les métriques runtime et l'historique.\\
Evaluateur pédagogique & Vérifier la conformité aux livrables MLOps/DataOps.\\
\bottomrule
\end{tabularx}
\end{table}

\section{Valeur Métier}

La valeur métier du projet repose sur quatre axes:

\begin{itemize}
    \item \textbf{Réactivité}: score de fraude disponible en temps quasi réel via \repoPath{POST /predict}.
    \item \textbf{Priorisation}: classement implicite des transactions selon la probabilité et le niveau de risque.
    \item \textbf{Traçabilité}: historique des prédictions, métriques, artefacts et versions de modèle.
    \item \textbf{Industrialisation}: pipeline reproductible, tests, CI/CD et déploiement containerisé.
\end{itemize}

\section{Data Strategy}

La stratégie data observée dans le dépôt est structurée autour d'un flux local contrôlé:

\begin{enumerate}
    \item \textbf{Source}: dataset Kaggle \repoPath{mlg-ulb/creditcardfraud}.
    \item \textbf{Cache local}: fichier \repoPath{data/raw/creditcard.csv}, non destiné à être versionné.
    \item \textbf{Ingestion}: chargement par dlt dans DuckDB.
    \item \textbf{Entrepôt local}: \repoPath{data/warehouse/fraud.duckdb}.
    \item \textbf{Transformation}: dbt construit les modèles staging et mart.
    \item \textbf{Contrat}: \repoPath{contracts/creditcard_fraud_contract.yml}.
    \item \textbf{Expérimentation}: MLflow pour paramètres, métriques, artefacts et registry.
    \item \textbf{Service}: FastAPI pour les prédictions et Streamlit pour l'interface.
\end{enumerate}
